\section{14.4 산업 사례: 자율주행 시뮬레이션 (Waymo의 Onboard/Offboard 루프)}
\begin{frame}{산업 규모로 확장: 자율주행 시뮬레이션}
  \begin{itemize}
    \item 14.1: XML로 로봇 정의 $\to$ ``환경이 바뀌어도 알고리즘은
      그대로'' 확인
    \item 14.2: Isaac Sim GPU 병렬의 ``수천 배''가 수학적으로
      어디서 오는지 보임
    \item 14.3: 정밀도 $\times$ 속도 두 축으로 도구 고르는 기준
  \end{itemize}
  \vskip0.6em
  \begin{block}{이번 절}
    ``\textbf{시뮬레이터를 실전 연구 수준으로 쓴다}''는 주제를
    \textbf{산업 규모, 즉 자율주행 시뮬레이션 스택}으로 확장.
    \kb{Waymo}의 자율주행 스택 구조를 예로:
    \begin{enumerate}
      \item 실차를 굴리는 \kb{Onboard}와 \kb{Offboard} 두 루프
      \item 이 구조가 강화학습의 \kb{policy iteration/Dyna}와 같은
        모양인 이유
      \item Offboard 루프를 현실적 시나리오로 채우는 diffusion
        기반 도구 \kb{SceneDiffuser}(Waymo, 2024)
    \end{enumerate}
  \end{block}
  {\footnotesize ML1 13.4와 15.4처럼 ``지도'' 수준 --- 깊은 유도 없이
  개념과 대표 논문을 잇는 가벼운 절.
  }
\end{frame}
